\documentclass[aspectratio=169,11pt]{beamer}

\usetheme{Madrid}
\setbeamertemplate{navigation symbols}{}
\setbeamertemplate{footline}[frame number]

\usepackage{fontspec}
\usepackage{polyglossia}
\setdefaultlanguage{russian}
\setotherlanguage{english}
\setsansfont{Arial}
\setmainfont{Times New Roman}
\setmonofont{Courier New}
\newfontfamily\cyrillicfont{Times New Roman}
\newfontfamily\cyrillicfontsf{Arial}
\newfontfamily\cyrillicfonttt{Courier New}

\usepackage{graphicx}
\usepackage{booktabs}
\usepackage{tabularx}
\usepackage{amsmath}
\usepackage{tikz}
\usetikzlibrary{arrows.meta, positioning, shapes.geometric}

\graphicspath{{slides/assets/}{thesis/figures/}}

\definecolor{phiBlue}{HTML}{0D3B66}
\definecolor{phiTeal}{HTML}{0F6E8C}
\definecolor{phiGreen}{HTML}{6AA84F}
\definecolor{phiGray}{HTML}{4D4D4D}

\newcommand{\tcell}[1]{\parbox[t]{\linewidth}{\raggedright\strut #1\strut}}

\setbeamercolor{title}{fg=white}
\setbeamercolor{frametitle}{fg=white, bg=phiBlue}
\setbeamercolor{block title}{fg=white, bg=phiTeal}
\setbeamercolor{block body}{fg=black, bg=black!3}
\setbeamercolor{item}{fg=phiBlue}
\setbeamercolor{structure}{fg=phiBlue}

\title{Исследование методов обработки данных\\в фазочувствительной рефлектометрии\\на длинных оптических линиях с усилителями}
\author{Студент группы Б02-212кт \\ Победин Николай Константинович \\[0.5ex]
    {\small Научный руководитель: Стрижак Даниил Валерьевич}}
\institute{\footnotesize Московский физико-технический институт (государственный университет) \\
    Физтех-школа физики и исследований им. Ландау \\
    Физтех-кластер академической и научной карьеры \\
    (Фундаментальные проблемы физики квантовых технологий)}
\date{2026}

\begin{document}

\begin{frame}
    \titlepage
\end{frame}

\begin{frame}[t]{Введение и обзор литературы}
    \footnotesize
    \begin{itemize}
        \item Оптическая рефлектометрия восстанавливает распределение обратного сигнала вдоль волокна по времени прихода рассеянного излучения.
        \item Физическая основа сигнала - \textbf{обратное Рэлеевское рассеяние} на случайных микронеоднородностях стекла.
        \item В классическом OTDR обычно анализируется усреднённый профиль мощности, поэтому для повышения отношения сигнал/шум требуется накопление большого числа трасс.
        \item В \textbf{phase-sensitive OTDR} используется когерентное зондирование: вклады обратного Рэлеевского рассеяния интерферируют, а внешнее воздействие меняет их относительные фазы.
        \item Поэтому информативны изменения уже \textbf{между соседними рефлектограммами}; на практике удобно работать с короткими окнами из десятков-сотен трасс, а не только с усреднённым профилем линии.
    \end{itemize}

    \vspace{0.2ex}
    \begin{alertblock}{Сложность задачи на длинных линиях с усилителями}
        \small
        На линиях в десятки километров затухание компенсируют EDFA-усилителями, но их ASE-шум делает фон \textbf{резко неоднородным по дистанции}, и классические энергетические детекторы теряют точность. Нужен автоматический алгоритм: парсинг, выравнивание, подавление фона и локализация возмущений.
    \end{alertblock}
\end{frame}

\begin{frame}{Цель и задачи исследования}
    \begin{block}{Цель}
        Разработать воспроизводимый метод обработки данных phi-OTDR, который позволяет автоматически выделять рефлектограммы из непрерывного сигнала, повышать интерпретируемость измерений и локализовать возмущения вдоль волоконной линии с распределёнными усилителями.
    \end{block}

    \vspace{0.5ex}
    \begin{itemize}
        \item Построить парсер, который сам определяет начала и длину рефлектограмм без внешней разметки (автооценка периода следования импульсов).
        \item Компенсировать межтрассовый джиттер кросс-корреляционным выравниванием с явными критериями приёма.
        \item Построить устойчивый к выбросам детектор: медианный фон, MAD-нормировка по дистанции и комбинированная спектральная оценка $S(g)$.
        \item Количественно валидировать метод на размеченном наборе из 130 файлов и сравнить с одно-компонентными детекторами-эталонами.
    \end{itemize}
\end{frame}

\begin{frame}[t]{Методы и материалы}
    \vspace{0.2ex}
    \begin{columns}[T,onlytextwidth]
        \column{0.55\textwidth}
        \raggedright
        \includegraphics[width=\linewidth,height=0.37\textheight,keepaspectratio,page=1]{phi_otdr_scheme.pdf}

        \column{0.40\textwidth}
        \includegraphics[width=\linewidth,height=0.4\textheight,keepaspectratio]{single_reflectogram_example.png}
    \end{columns}

    \vspace{0.4ex}
    \noindent
    \begin{minipage}[t]{0.50\textwidth}
        \begin{block}{Экспериментальные данные}
            \scriptsize
            \begin{itemize}
                \item 882 записи (около 30 мин) с 4 линий, длина до ${\sim}60$ км;
                \item АЦП 50 МГц, $\lambda = 1550$ нм, импульс $\tau \in \{100, 200, 300, 500, 1000\}$ нс;
                \item размеченный поднабор: \textbf{130 файлов} (92 с воздействием, 38 без).
            \end{itemize}
        \end{block}
    \end{minipage}
    \hfill
    \begin{minipage}[t]{0.27\textwidth}
        \begin{block}{Физическая интерпретация}
            \scriptsize
            Координата вдоль линии вычисляется как
            \[
                x = \frac{c\,t}{2n}, \qquad n \approx 1.468.
            \]
        \end{block}
    \end{minipage}
\end{frame}

\begin{frame}[t]{Стратегия проверки: proof of concept и валидация}
    \footnotesize
    \begin{block}{Этап 1. Proof of concept - контролируемая запись}
        Разработка и демонстрация пайплайна на \textbf{контролируемом модельном воздействии} (гармоника 50 Гц) на \textbf{длинной} реальной линии ($\sim$60 км) с распределённым EDFA-усилением. Цель - показать, что детекция:
        \begin{itemize}
            \setlength{\itemsep}{1pt}
            \item уверенно локализует возмущение \textbf{на длинной линии} (водопад и компоненты $S(g)$ - на следующих слайдах);
            \item работает на \textbf{очень коротком окне} наблюдения $\sim$0,1 с (100 рефлектограмм).
        \end{itemize}
    \end{block}
    \begin{block}{Этап 2. Валидация подхода - данные, приближенные к реальности}
        Количественная оценка - со стенда с \textbf{имитацией шагов по грунту} (реалистичный сценарий вторжения) на \textbf{длинной линии $\sim$500 км} (10-й пролёт, EDFA). Анализ ведётся в \textbf{окне $\sim$6-9 км вокруг события} (всю трассу хранить не нужно). Медианная ошибка локализации \textbf{61 м}, F1 = 0,48 на 130 файлах.
    \end{block}
\end{frame}

\begin{frame}[t]{Выровненные рефлектограммы и спектральный baseline}
    \begin{columns}[T,onlytextwidth]
        \column{0.50\textwidth}
        \includegraphics[width=\linewidth,height=0.54\textheight,keepaspectratio]{aligned_waterfall_presentation.png}

        \column{0.50\textwidth}
        \includegraphics[width=\linewidth,height=0.54\textheight,keepaspectratio]{fourier_full_range_presentation.png}
    \end{columns}

    \vspace{0.5ex}
    \scriptsize
    \begin{itemize}
        \item После парсинга, кросс-корреляционного выравнивания и перевода в координаты расстояние - время видно, что вдоль линии существуют локальные аномальные зоны.
        \item Fourier/STFT-карта подтверждает наличие спектрального следа возмущения и его локализацию по расстоянию.
        \item Но полный спектральный просмотр для всех точек линии дорог: для окна из $K$ рефлектограмм длиной $M$ отсчётов лобовая карта без группировки масштабируется как $O(M\,K\log K)$, поэтому дальше спектральная информация сворачивается в компактный автоматический score.
    \end{itemize}
\end{frame}

\begin{frame}[t]{Алгоритм детекции: идея и математическая схема}
    \begin{center}
        \resizebox{0.96\textwidth}{!}{
            \begin{tikzpicture}[
                node distance=0.55cm,
                block/.style={draw=phiBlue, rounded corners=2pt, fill=phiBlue!7, align=center, minimum width=2.3cm, minimum height=0.95cm},
                arrow/.style={-{Latex[length=2mm]}, thick, draw=phiBlue}
            ]
                \node[block] (a) {Aligned\\reflectograms};
                \node[block, right=of a] (b) {Median\\background};
                \node[block, right=of b] (c) {Residual\\signal};
                \node[block, right=of c] (d) {MAD\\normalization};
                \node[block, right=of d] (e) {FFT-based\\features};
                \node[block, right=of e] (f) {Combined\\score + peaks};

                \draw[arrow] (a) -- (b);
                \draw[arrow] (b) -- (c);
                \draw[arrow] (c) -- (d);
                \draw[arrow] (d) -- (e);
                \draw[arrow] (e) -- (f);
            \end{tikzpicture}
        }
    \end{center}

    \vspace{0.2ex}
    \begin{block}{Робастная предобработка}
        \small
        \[
            b(x)=\operatorname{med}_t I(x,t), \qquad
            r(x,t)=I(x,t)-b(x)
        \]
        \[
            s(x)=1.4826 \cdot \operatorname{med}_t \left| r(x,t)-\operatorname{med}_t r(x,t)\right|, \qquad
            z(x,t)=\frac{r(x,t)}{s(x)+\varepsilon}
        \]
    \end{block}

    \small
    \begin{itemize}
        \item $b(x)$ выделяет \textbf{стационарный профиль линии}: затухание, усилители, устойчивые отражающие участки.
        \item $r(x,t)$ убирает постоянный фон и оставляет только отклонения относительно типичной рефлектограммы.
        \item $s(x)$ задаёт \textbf{робастный масштаб} по расстоянию, а $z(x,t)$ становится входом для частотных и энергетических признаков.
    \end{itemize}
\end{frame}

\begin{frame}[t]{Алгоритм детекции: фон и нормировка}
    \begin{columns}[T,onlytextwidth]
        \column{0.64\textwidth}
        \includegraphics[width=\linewidth,height=0.39\textheight,keepaspectratio]{detect_step_background.png}

        \vspace{0.4ex}
        \includegraphics[width=\linewidth,height=0.39\textheight,keepaspectratio]{detect_step_mad_profile.png}

        \column{0.36\textwidth}
        {\footnotesize
        \textbf{Что происходит на этом этапе}
        \begin{itemize}
            \item Медианный фон $b(x)$ описывает устойчивую форму линии как функцию расстояния.
            \item MAD-масштаб показывает, где естественный разброс по времени велик, а где мал.
            \item После нормировки сравнивается не абсолютная яркость, а статистически значимое отклонение.
        \end{itemize}

        \vspace{0.2ex}
        \textbf{Робастная нормировка}
        \[
            z(x,t)=\frac{I(x,t)-\operatorname{med}_t I(x,t)}{1.4826\operatorname{MAD}_t(r)+\varepsilon}.
        \]
        }
    \end{columns}
\end{frame}

\begin{frame}[t]{Алгоритм детекции: отдельные признаки}
    \begin{center}
        \includegraphics[width=0.97\linewidth,height=0.62\textheight,keepaspectratio]{detect_score_components_2x2_wide.png}
    \end{center}

    \scriptsize
    \begin{itemize}
        \item $S_{\text{peak}}$ реагирует на узкополосные компоненты, $S_{\text{broad}}$ - на широкополосные спектральные отклонения.
        \item $E_z$ страхует случаи, когда событие видно по энергии во времени, но не доминирует на одной частоте; итоговый $S(g)=0{,}55\,S_{\text{peak}}+0{,}30\,S_{\text{broad}}+0{,}15\,E_z$ строится как их взвешенная сумма.
    \end{itemize}
\end{frame}

\begin{frame}[t]{Короткое окно: 100 рефлектограмм}
    \begin{columns}[T,onlytextwidth]
        \column{0.70\textwidth}
        \includegraphics[width=\linewidth,height=0.44\textheight,keepaspectratio]{short_window_score_100_presentation.png}

        \column{0.30\textwidth}
        {\scriptsize
        \textbf{Что показывает окно 100 трасс}
        \begin{itemize}
            \item При $f_{\mathrm{rep}}\approx 1000$ Гц это всего \textbf{0.1 с} наблюдения.
            \item Обе размеченные зоны (около \textbf{32 км} и \textbf{42 км}) уверенно выделяются: пики $S(g)$ более чем в \textbf{2.5 раза} превышают порог, ложных срабатываний нет.
            \item На полной односекундной записи контраст сигнал/фон практически тот же.
        \end{itemize}
        }
    \end{columns}

    \vspace{0.5ex}
    {\scriptsize
    \textbf{Преимущества предлагаемого подхода}
    \begin{itemize}
        \item Спектральная информация всё ещё извлекается из всей записи, но в качестве конечного результата сохраняется не полная двумерная карта, а только \textbf{одномерный score по расстоянию}.
        \item Это упрощает хранение, последующий анализ и автоматическое принятие решения: вместо time-frequency heatmap нужно рассматривать только один профиль вдоль линии.
        \item В отличие от классического OTDR здесь сохраняется фазовая интерференционная динамика соседних трасс, поэтому информативны уже короткие окна измерения - без отдельной эталонной записи "спокойного" состояния линии.
    \end{itemize}
    }
\end{frame}

\begin{frame}[t]{Итоговая локализация возмущений}
    \begin{columns}[T,onlytextwidth]
        \column{0.50\textwidth}
        \includegraphics[width=\linewidth,height=0.56\textheight,keepaspectratio]{waterfall_zoom_30_35_presentation.png}

        \column{0.50\textwidth}
        \includegraphics[width=\linewidth,height=0.56\textheight,keepaspectratio]{waterfall_zoom_40_45_presentation.png}
    \end{columns}

    \vspace{0.7ex}
    \small
    \textbf{Контрольный случай (proof of concept, длинная линия):} две зоны возмущения (около 32 км и 42 км) локализованы как \textbf{31.7 км} и \textbf{42.6 км}. Количественная точность - на реалистичном наборе (следующий слайд): медианная ошибка локализации \textbf{61 м} (в окне анализа $\sim$6-9 км вокруг события).
\end{frame}

\begin{frame}[t]{Валидация на реалистичных данных (имитация шагов)}
    \footnotesize
    \begin{columns}[T,onlytextwidth]
        \column{0.46\textwidth}
        \begin{block}{Протокол оценки}
            \scriptsize
            \begin{itemize}
                \item реалистичный сценарий: стенд с имитацией шагов по грунту, длинная линия $\sim$500 км (10-й пролёт), анализ в окне $\sim$6-9 км вокруг события;
                \item 130 файлов: 92 с воздействием, 38 без;
                \item единый набор гиперпараметров ($G=20$, $f_\text{min}=5$ Гц, $\Delta_\text{min}=2$ км);
                \item попадание - детекция в коридоре $\pm 0.5$ км вокруг размеченной зоны, всё остальное - FP.
            \end{itemize}
        \end{block}

        \column{0.50\textwidth}
        \vspace{0.5ex}
        \renewcommand{\arraystretch}{1.3}
        \begin{tabular}{@{}l r@{}}
            \toprule
            \textbf{Метрика} & \textbf{Значение} \\
            \midrule
            Recall (порог) & 0{,}38 \;(35/92) \\
            Precision (порог) & 0{,}64 \;(FP $=20$) \\
            \textbf{F1-мера} & \textbf{0{,}48} \\
            Recall (топ-кандидат) & 0{,}47 \;(43/92) \\
            Медианная ошибка локализации & 61 м \\
            \bottomrule
        \end{tabular}
    \end{columns}

    \vspace{1.0ex}
    \scriptsize
    \begin{itemize}
        \item 31 из 35 попаданий локализованы с ошибкой $\le 0.2$ км - на уровне пространственного разрешения установки.
        \item Recall выше всего на средних импульсах 200-500 нс ($0{,}50-0{,}59$); на 100 нс мешает низкий SNR одиночного отсчёта, на 1000 нс - пространственное размытие события на ${\sim}100$ м.
    \end{itemize}
\end{frame}

\begin{frame}[t]{Зачем спектральные признаки: сравнение с эталонами}
    \begin{columns}[T,onlytextwidth]
        \column{0.52\textwidth}
        \includegraphics[width=\linewidth,height=0.55\textheight,keepaspectratio]{fig_f1_bars.png}

        \vspace{1.2ex}
        \begin{block}{Главный вывод}
            \scriptsize
            Лучший баланс Recall/Precision даёт только комбинированная оценка $S(g)$ с двойным порогом. Одиночные признаки (пик, энергия) набирают меньший F1 и кратно больше ложных срабатываний на неоднородном фоне линии.
        \end{block}

        \column{0.46\textwidth}
        \footnotesize
        Все варианты используют один и тот же парсер, выравнивание и предобработку; отличается только источник итоговой оценки.

        \vspace{0.6ex}
        \renewcommand{\arraystretch}{1.25}
        \setlength{\tabcolsep}{4pt}
        \scriptsize
        \begin{tabular}{@{}l c c c@{}}
            \toprule
            \textbf{Детектор} & \textbf{F1} & \textbf{Prec.} & \textbf{FP} \\
            \midrule
            Комбинированный & \textbf{0{,}48} & \textbf{0{,}64} & 20 \\
            Только пик & 0{,}43 & 0{,}38 & 71 \\
            Только энергия & 0{,}33 & 0{,}29 & 89 \\
            \bottomrule
        \end{tabular}

        \vspace{0.8ex}
        \begin{itemize}
            \item Энергетический детектор - слабейший: на неоднородном фоне он не отличает локальное событие от ASE-шума вблизи усилителей (в 4.5 раза больше FP).
            \item Комбинация спектральных признаков с MAD-нормировкой по дистанции даёт прирост F1 не менее чем на 5 п.п.
        \end{itemize}
    \end{columns}
\end{frame}

\begin{frame}[t]{Сравнение с подходами из литературы}
    \fontsize{6.2pt}{6.9pt}\selectfont
    \renewcommand{\arraystretch}{1.05}
    \setlength{\tabcolsep}{4pt}
    \begin{center}
    \begin{tabular}{@{}>{\raggedright\arraybackslash}p{0.20\textwidth}>{\raggedright\arraybackslash}p{0.22\textwidth}>{\raggedright\arraybackslash}p{0.22\textwidth}>{\raggedright\arraybackslash}p{0.27\textwidth}@{}}
        \toprule
        \textbf{Метод} & \textbf{Применимость} & \textbf{Выход} & \textbf{Ограничения} \\
        \midrule
        \tcell{CWT (Qin, 2012)} &
        \tcell{Короткие линии, нестационарные импульсы} &
        \tcell{2D скалограмма + гребень} &
        \tcell{Полный 2D-анализ, нет учёта неоднородного фона} \\
        \midrule
        \tcell{HHT (Hui, 2014)} &
        \tcell{Нестационарные сигналы, адаптивный базис} &
        \tcell{2D распределение (мода, дистанция)} &
        \tcell{Артефакты перемежения мод, чувствительность к шуму} \\
        \midrule
        \tcell{Вейвлет-пакеты (Li, 2014)} &
        \tcell{Стационарные вибрации} &
        \tcell{2D карта (пакет, дистанция)} &
        \tcell{Двумерный выход, требует дополнительной агрегации} \\
        \midrule
        \tcell{Сопоставление матриц (Liu, 2018)} &
        \tcell{Короткие линии (1-2 км), real-time} &
        \tcell{Бинарная карта расстояний} &
        \tcell{Зависимость от эталонного состояния; короткие линии} \\
        \midrule
        \tcell{Машинное обучение (HMM, CNN)} &
        \tcell{Классификация типа события} &
        \tcell{Метка класса} &
        \tcell{Нужна большая разметка, нет физической интерпретируемости} \\
        \midrule
        \tcell{\textbf{Предлагаемый робастный детектор}} &
        \tcell{\textbf{Длинные линии с неоднородным фоном (EDFA)}} &
        \tcell{\textbf{1D score $S(g)$ по дистанции}} &
        \tcell{\textbf{Чувствителен к качеству парсинга и калибровке порогов}} \\
        \bottomrule
    \end{tabular}
    \end{center}

    \vspace{0.6ex}
    \scriptsize
    Общий пробел литературы: методы либо ориентированы на короткие линии и не учитывают неоднородность фона, либо выдают 2D-карту в плоскости (частота, дистанция), либо требуют большой разметки. Предлагаемый метод закрывает именно этот пробел - компактный одномерный результат на длинной линии с усилителями.
\end{frame}

\begin{frame}{Выводы}
    \begin{itemize}
        \item Разработан и валидирован робастный алгоритм детекции возмущений в данных phi-OTDR для длинных волоконных линий с EDFA-усилителями и неоднородным фоном по расстоянию.
        \item Ключевое преимущество - свёртка time-frequency информации в \textbf{одномерный score $S(g)$ по расстоянию}, который существенно проще хранить и интерпретировать, чем двумерные спектральные карты.
        \item На размеченном наборе из 130 файлов достигнуты \textbf{F1\,=\,0{,}48} (Precision 0{,}64, Recall 0{,}38), медианная ошибка локализации \textbf{61 м}.
        \item Спектральные признаки критичны: энергетический детектор-эталон даёт лишь F1\,=\,0{,}33 - на неоднородном фоне он тонет в ASE-шуме усилителей.
        \item Метод сохраняет чувствительность уже на \textbf{коротких окнах ${\sim}0.1$ с} (порядка 100 рефлектограмм) и не требует эталонной записи "спокойного" состояния линии.
        \item Дальнейшая работа: скользящая нормировка фона на длинных (${\sim}10$ с) записях, адаптивный размер группы $G$ под длинный импульс, расширение отрицательного набора (текущий Precision - оценка сверху).
    \end{itemize}
\end{frame}

% =====================================================================
%  РЕЗЕРВНЫЕ СЛАЙДЫ (для вопросов)
% =====================================================================
\appendix

\begin{frame}[t]{Резервный слайд: оптимальность весов в $S(g)$}
    \begin{columns}[T,onlytextwidth]
        \column{0.42\textwidth}
        \includegraphics[width=\linewidth,height=0.47\textheight,keepaspectratio]{weight_sweep_ternary.png}

        \column{0.58\textwidth}
        \vspace{1.2ex}
        \includegraphics[width=\linewidth,height=0.43\textheight,keepaspectratio]{weight_sweep_hist.png}
    \end{columns}

    \vspace{0.6ex}
    \scriptsize
    \begin{itemize}
        \setlength{\itemsep}{1.2pt}
        \item Варьируются \textbf{только веса} $w_1,w_2,w_3$ в $S(g)=w_1\,peak+w_2\,broad+w_3\,energy$ ($w_i\ge 0$, $\sum w_i=1$); остальная машина детектора (двойной порог, peak-guard) фиксирована - это контролируемый эксперимент именно про коэффициенты.
        \item Выбранные $(0{,}55;\,0{,}30;\,0{,}15)$: \textbf{F1\,=\,0{,}48} - выше среднего по случайным комбинациям ($0{,}46$) и любого вырожденного набора «всё на один компонент» (peak $0{,}46$, broad $0{,}44$, energy $0{,}43$).
        \item На высоком плато: бьёт \textbf{83\,\%} случайных наборов; до максимума на сетке ($0{,}49$) - $1{,}4$ п.п. F1 ($\approx$1-2 события из 92). Веса зафиксированы на пилотных записях, \textbf{не} подгонялись под тестовый набор.
    \end{itemize}
\end{frame}

\begin{frame}[t]{Резервный слайд: типы воздействий и длина линии}
    \centering
    \includegraphics[width=0.99\linewidth,height=0.56\textheight,keepaspectratio]{dataset_composition.png}

    \vspace{0.6ex}
    \footnotesize
    \begin{block}{Важно: на каких линиях измерено}
        \scriptsize
        Полный датасет ($\approx$882 записи, 4 линии) содержит разные воздействия, но количественная валидация (F1\,=\,0,48 на 130 файлах) проведена \textbf{только на шагах} - длинная линия $\sim$500 км (10-й пролёт), анализ в окне $\sim$6-9 км вокруг события. Гармоника, фон и растяжка сняты на линии до $\sim$60 км (окно регистрации $\sim$100 км), нагрев - на \textbf{реальных длинных линиях 50-950 км} (шкаф катушек, усилители EDFA между пролётами); эти типы \textbf{не размечены} и в метрики не входят.
    \end{block}
\end{frame}

\end{document}
